% ===== Preambul comun pentru examen, banca de probleme și setul de exerciții (XeLaTeX) =====
% Statistica piețelor financiare / Statistics of Financial Markets -- Daniel Traian Pele
% Folosire: \documentclass[10.5pt]{article}  \def\examlang{ro} (sau en)  \input{../sfm_exam_preamble}
% Fișierele se compilează din exam/<subfolder>/, deci logo-urile sînt în ../../logos/.
\usepackage{fontspec}
\setmainfont{Helvetica Neue}
\usepackage{amsmath,amssymb}
\usepackage[a4paper,margin=1.9cm,top=2.4cm,headsep=0.5cm,headheight=14pt]{geometry}
\usepackage{graphicx}
\usepackage{float}
\usepackage{booktabs}
\usepackage{array}
\usepackage{enumitem}
\usepackage{xcolor}
\usepackage{fancyhdr}
\usepackage{lastpage}
\usepackage[hidelinks]{hyperref}
\definecolor{spfblue}{RGB}{26,58,110}
\definecolor{spfred}{RGB}{205,0,0}
\graphicspath{{../../logos/}{figs/}{../bank/figs/}{../practice/figs/}}

\providecommand{\examlang}{ro}
\newif\ifen
\def\tmpen{en}\ifx\examlang\tmpen\entrue\else\enfalse\fi

% etichete
\ifen
  \renewcommand{\figurename}{Figure}\renewcommand{\tablename}{Table}
  \newcommand{\coursetitle}{Statistics of Financial Markets}
  \newcommand{\subjword}{Problem}
  \newcommand{\rezword}{Solution.}
  \newcommand{\baremword}{Grading:}
\else
  \renewcommand{\figurename}{Figura}\renewcommand{\tablename}{Tabelul}
  \newcommand{\coursetitle}{Statistica piețelor financiare}
  \newcommand{\subjword}{Subiectul}
  \newcommand{\rezword}{Rezolvare.}
  \newcommand{\baremword}{Barem:}
\fi

% comutator: barem/rezolvare (true) sau subiecte (false)
\newif\ifrez \rezfalse

% antet / subsol
\pagestyle{fancy}
\fancyhf{}
\renewcommand{\headrulewidth}{0.4pt}
\lhead{\small\itshape \coursetitle}
\providecommand{\headyear}{2026}
\rhead{\small\bfseries \headyear}
\cfoot{\small\thepage/\pageref{LastPage}}

\setlength{\parindent}{0pt}
\setlength{\parskip}{0.42ex}
\setlength{\emergencystretch}{3em}

% numerotare automată, continuă, a cerințelor (1, 2, 3, ...)
\newcounter{qq}
\newcommand{\Q}{\refstepcounter{qq}\textbf{\theqq.}~}
% punctaj aliniat la dreapta, pe ultima linie a cerinței
\newcommand{\pts}[1]{\nobreak\hfill\textnormal{\itshape (#1)}\par}

% titlu de subiect: "Subiectul I • Titlu ........ puncte" + linie
\newcommand{\subiect}[3]{%
  \vspace{1.05ex}%
  {\color{spfblue}\bfseries \subjword\ #1 \textbullet\ #2}\hfill{\bfseries #3}\par
  \vspace{0.2ex}{\color{spfblue}\hrule height0.8pt}\vspace{0.7ex}}
% subiect numerotat automat (I, II, ...); argumentul opțional = identificatorul din bancă (vizibil doar în barem)
\newcounter{subj}
\newcommand{\problem}[3][]{%
  \stepcounter{subj}%
  \subiect{\Roman{subj}}{#2\ifrez\if\relax\detokenize{#1}\relax\else\ {\normalfont\footnotesize\color{spfred}[#1]}\fi\fi}{#3}}

% blocul de rezolvare (doar în barem)
\newcommand{\Rez}{\par\vspace{0.2ex}\textit{\rezword}\par\nopagebreak}
\newcommand{\barem}[1]{\par\textit{\baremword\ #1}\par}

% bloc de titlu cu sigle
\newcommand{\titlublock}[1]{%
\begin{center}
  \begin{minipage}[c]{0.16\textwidth}\includegraphics[height=1.15cm]{ase_logo.png}\end{minipage}%
  \begin{minipage}[c]{0.68\textwidth}\centering
  \ifen
    {\bfseries Bucharest University of Economic Studies}\\[0.2ex]
    Faculty of Cybernetics, Statistics and Economic Informatics\\[0.2ex]
    Course: Statistics of Financial Markets
  \else
    {\bfseries Academia de Studii Economice din București}\\[0.2ex]
    Facultatea de Cibernetică, Statistică și Informatică Economică\\[0.2ex]
    Disciplina: Statistica piețelor financiare
  \fi
  \end{minipage}%
  \begin{minipage}[c]{0.16\textwidth}\raggedleft\includegraphics[height=0.95cm]{ida_logo.png}\end{minipage}\\[1.2ex]
  {\large\bfseries #1}
\end{center}}

% valori utile generale (z_alpha = cuantila de ordin alpha a distribuției Normale standard)
\newcommand{\valoriutilegen}{%
\ifen Useful values: $z_{0.05}=-1.645$, $z_{0.025}=-1.960$, $z_{0.01}=-2.326$, $\varphi(1.960)=0.0584$,
$\varphi(2.326)=0.0267$, $\ln 2=0.693$, $\sqrt{252}\approx 15.87$, $\sqrt{365}\approx 19.10$,
$\chi^2_{0.95}(1)=3.84$, $\chi^2_{0.95}(2)=5.99$, $\chi^2_{0.95}(10)=18.31$
($z_p$ and $\chi^2_p$ are quantiles of order $p$; $\varphi$ is the standard Normal density).
\else Valori utile: $z_{0,05}=-1{,}645$; $z_{0,025}=-1{,}960$; $z_{0,01}=-2{,}326$; $\varphi(1{,}960)=0{,}0584$;
$\varphi(2{,}326)=0{,}0267$; $\ln 2=0{,}693$; $\sqrt{252}\approx 15{,}87$; $\sqrt{365}\approx 19{,}10$;
$\chi^2_{0,95}(1)=3{,}84$; $\chi^2_{0,95}(2)=5{,}99$; $\chi^2_{0,95}(10)=18{,}31$
($z_p$ și $\chi^2_p$ sînt cuantile de ordin $p$; $\varphi$ este densitatea distribuției Normale standard).\fi}

% titlul capitolului într-un fișier din bancă (folosit doar ca etichetă de build_variant.py)
\newcommand{\banktitle}[2]{}
